% !TeX spellcheck = pt_BR
\documentclass[12pt,portuguese]{article}

%----------------------------------------------------------------------
% PACOTES ESSENCIAIS E TRADUÇÃO
%----------------------------------------------------------------------
\usepackage[brazilian]{babel}     % Suporte à língua portuguesa
\usepackage[utf8]{inputenc}       % Codificação de entrada
\usepackage[T1]{fontenc}          % Codificação da fonte
\usepackage{lmodern}              % Fonte escalável

%----------------------------------------------------------------------
% PACOTES MATEMÁTICOS
%----------------------------------------------------------------------
\usepackage{amsmath, amsthm, amssymb, amstext, mathtools}
\usepackage{mathrsfs}             % Para fontes como \mathscr
\usepackage{dsfont}               % Para o indicador de conjunto \mathds{1}
\usepackage[all]{xy}              % Para diagramas comutativos
\usepackage{imakeidx}             % Criação de Lista de Símbolos / Índice

% Desigualdades sempre com símbolos inclinados
\let\leq\leqslant
\let\le\leqslant
\let\ge\geqslant
\let\geq\geqslant

%----------------------------------------------------------------------
% FORMATAÇÃO, LAYOUT E CORES
%----------------------------------------------------------------------
\usepackage{geometry}
\usepackage{varwidth}
\usepackage{graphicx}
\usepackage[pdftex]{xcolor}
\usepackage{titlesec}
\usepackage{appendix}
\usepackage{enumitem}
\usepackage{float}
\usepackage{booktabs}
\usepackage{array}
\usepackage{colortbl}
\usepackage{fontawesome5}
\usepackage{lucide-icons}
\usepackage{caption}
\usepackage{stackengine}
\captionsetup{width=0.8\textwidth}

% Definição de Cores
\definecolor{teal}{RGB}{15,118,110}
\definecolor{text}{RGB}{28,25,23}
\definecolor{muted}{RGB}{87,83,78}
\definecolor{card}{RGB}{255,255,255}
\definecolor{border}{RGB}{231,229,228}

% Definição de margens e espaçamento
\setlength{\topmargin}{-0.7in}
\setlength{\oddsidemargin}{0in}
\setlength{\textheight}{9.55in}
\setlength{\textwidth}{6.5in}
\setlength{\baselineskip}{12mm}

% Formatação dos títulos de seção
\titleformat{\section}{\Large\bfseries}{\thesection.}{0.5em}{}
\titleformat{\subsection}{\large\bfseries}{\thesubsection.}{0.5em}{}
\titleformat{\subsubsection}{\normalfont\itshape}{\thesubsubsection.}{0.6em}{}

%----------------------------------------------------------------------
% TIKZ E GRÁFICOS
%----------------------------------------------------------------------
\usepackage{tikz}
\usepackage{pgfplots}
\pgfplotsset{compat=1.17}
\usetikzlibrary{positioning,calc,fit,nfold,decorations.pathreplacing,calligraphy,tikzmark,patterns,arrows.meta,bending,hobby}

%----------------------------------------------------------------------
% REFERÊNCIAS, LINKS E METADADOS DO PDF
%----------------------------------------------------------------------
\usepackage[pdftex, colorlinks, linkcolor={red!50!black}, 
citecolor={red}, urlcolor={blue!80!black}]{hyperref}

\hypersetup{
pdftitle={Funções de Variáveis Aleatórias Contínuas e Esperança},
pdfauthor={L. Cioletti},
pdfsubject={Funções de variáveis aleatórias contínuas, transformações de densidades, esperança, variância, lei do estatístico subconsciente},
pdfkeywords={variáveis aleatórias contínuas, função densidade, transformação de variáveis, esperança, variância, MAT-UnB}
}

% --- PACOTE E COMANDO PARA A LICENÇA CREATIVE COMMONS ---
\usepackage[type={CC},modifier={by-nc-sa},version={4.0}]{doclicense}
\newcommand\blfootnote[1]{%
\begingroup
\renewcommand\thefootnote{}\footnote{#1}%
\addtocounter{footnote}{-1}%
\endgroup
}

%----------------------------------------------------------------------
% ALIASCNT E AMBIENTES DE TEOREMA (Cleveref)
%----------------------------------------------------------------------
\usepackage{aliascnt}
\theoremstyle{definition}
\newtheorem{theorem}{Teorema}

\newaliascnt{definition}{theorem}
\newtheorem{definition}[definition]{Definição}
\aliascntresetthe{definition}

\newaliascnt{example}{theorem}
\newtheorem{example}[example]{Exemplo}
\aliascntresetthe{example}

\newaliascnt{proposition}{theorem}
\newtheorem{proposition}[proposition]{Proposição}
\aliascntresetthe{proposition}

\newaliascnt{corollary}{theorem}
\newtheorem{corollary}[corollary]{Corolário}
\aliascntresetthe{corollary}

\newaliascnt{lemma}{theorem}
\newtheorem{lemma}[lemma]{Lema}
\aliascntresetthe{lemma}

\newaliascnt{remark}{theorem}
\newtheorem{remark}[remark]{Observação}
\aliascntresetthe{remark}

\newtheoremstyle{plainexercise}
  {3pt}{3pt}{\normalfont}{}{\bfseries}{.}{.5em}
  {\thmnumber{#2}\thmnote{ #3}}
\theoremstyle{plainexercise}
\newtheorem{exercise}{Exercício}

% Versão marcada dos exercícios: a estrela fica pendurada na margem esquerda,
% como nas notas anteriores.
\newtheoremstyle{starredexercise}
  {3pt}{3pt}{\normalfont}{}{\bfseries}{.}{.5em}
  {\llap{\textasteriskcentered\,}\thmnumber{#2}\thmnote{ #3}}
\theoremstyle{starredexercise}
\newtheorem{starredexercise}[exercise]{Exercício}
\theoremstyle{definition}

% Ambientes não numerados usados no texto
\newtheorem*{property}{Propriedade}
\newtheorem*{solution}{Solução}

% Se um enumerate for o primeiro conteúdo de um teorema/exercício (cabeçalho
% ainda pendente, \@inlabel true), o \leavevmode dispara o \everypar do amsthm
% que coloca o cabeçalho na própria linha, antes da lista iniciar.
\makeatletter
\setlist[enumerate]{before={\if@inlabel\leavevmode\fi}}
\makeatother

% Nomes em português para os ambientes
\AtBeginDocument{\renewcommand{\proofname}{\normalfont\textbf{Prova}}}
\renewcommand{\qedsymbol}{$\blacksquare$}

% Cleveref (deve vir depois de hyperref e amsthm)
\usepackage[brazilian,nameinlink]{cleveref}

\crefname{theorem}{teorema}{teoremas}
\Crefname{theorem}{Teorema}{Teoremas}
\crefname{definition}{definição}{definições}
\Crefname{definition}{Definição}{Definições}
\crefname{lemma}{lema}{lemas}
\Crefname{lemma}{Lema}{Lemas}
\crefname{proposition}{proposição}{proposições}
\Crefname{proposition}{Proposição}{Proposições}
\crefname{corollary}{corolário}{corolários}
\Crefname{corollary}{Corolário}{Corolários}
\crefname{example}{exemplo}{exemplos}
\Crefname{example}{Exemplo}{Exemplos}
\crefname{exercise}{exercício}{exercícios}
\Crefname{exercise}{Exercício}{Exercícios}
\crefname{remark}{observação}{observações}
\Crefname{remark}{Observação}{Observações}
\crefname{section}{seção}{seções}
\Crefname{section}{Seção}{Seções}

%----------------------------------------------------------------------
% MACROS ADICIONAIS
%----------------------------------------------------------------------
\let\oldemptyset\emptyset
\let\emptyset\varnothing
\graphicspath{{./}{../arquivos-aux/}}

% Seno na convenção em português
\renewcommand{\sin}{\operatorname{sen}}

%----------------------------------------------------------------------
% DOCUMENTO PRINCIPAL
%----------------------------------------------------------------------
\begin{document}
% --- CABEÇALHO COM O LOGO DA UnB ---
\begin{figure}[H]
\includegraphics{logo-mat-unb.pdf}
\end{figure}

\vspace{0.1cm}

% --- TÍTULO, AUTOR E DATA ---
\begin{center}
{\LARGE\bf Funções de Variáveis Aleatórias Contínuas e Esperança}
\vspace{1cm}

{\large \bf L. Cioletti}
\vspace{0.3cm}

{\footnotesize Setembro de 2026}
\end{center}

% --- RESUMO (ABSTRACT) ---
\begin{abstract}
Nestas notas damos continuidade ao estudo das variáveis aleatórias iniciado nas
notas sobre funções de distribuição. Mostramos como determinar a função
densidade de uma variável aleatória da forma $Y=g(X)$ a partir da densidade
de $X$, quando a função $g$ é uma função monótona ou quadrática. Em seguida
definimos a esperança de uma variável aleatória absolutamente contínua e
apresentamos o enunciado do Teorema da Mudança de Variáveis, que
permite, em particular, calcular $\mathbb{E}[g(X)]$ diretamente. 
Em seguida, introduzimos a variância como medida de dispersão, mostrando a identidade
$\operatorname{Var}(X)=\mathbb{E}[X^2]-\mathbb{E}[X]^2$. A exposição segue de
perto as Seções 5.5 e 5.6 do Capítulo 5 de \cite{GrimmettWelsh2014}. Como
aplicação, deduzimos a desigualdade de Cauchy-Schwarz a partir do fato de que
a variância é sempre não negativa e a traduzimos em uma desigualdade sobre
integrais. A nota conta com uma lista de exercícios ao final.
\end{abstract}

% --- LICENÇA NO RODAPÉ ---
\blfootnote{%
\doclicenseImage[imagewidth=2.6em]\ 
\begingroup
\hypersetup{urlcolor=darkgray}
\href{\doclicenseURL}{%
\doclicenseName - Attribution-NonCommercial-ShareAlike 4.0 International%
}%
\endgroup
}

%----------------------------------------------------------------------
% CONTEÚDO
%----------------------------------------------------------------------

\section{Funções de Variáveis Aleatórias Contínuas}\label{sec:funcoes-va-continuas}

Nas notas anteriores estudamos as variáveis aleatórias absolutamente contínuas
e suas funções densidade. O problema que estamos interessado agora é o seguinte.
Seja $X$ uma variável aleatória sobre um espaço de probabilidade 
$(\Omega, \mathcal{F}, \mathbb{P})$ e suponha que $g:\mathbb{R} \to \mathbb{R}$
seja uma função. Sob quais condições a composição $Y \equiv g(X)$ que é uma
aplicação de $\Omega$ em $\mathbb{R}$ definida em cada ponto $\omega\in\Omega$ 
pela expressão $Y(\omega) \equiv g(X(\omega))$ é uma variável aleatória?  Veja o diagrama abaixo.

\begin{figure}[H]
\centering
\begin{tikzpicture}[
  >=Stealth,
  ponto/.style={circle, fill=black, inner sep=0pt, minimum size=1.2mm},
  blob/.style={use Hobby shortcut, closed, line width=0.6pt},
]
  \draw[blob]
    (-1.19,0.9) .. (-0.08,-0.94) .. (0.78,-0.47) .. (1.9,-1.07) .. (2.71,0.1)
    .. (2.91,1.06) .. (2.4,1.59) .. (1.29,1.77) .. (0.17,1.91);
  \node at (0.8,2.35) {$\Omega$};
  \node[ponto] (omega) at (0.6,0.4) {};
  \node[below=1mm] at (0.6,0.4) {$\omega$};

  \draw[->, line width=0.5pt] (7.4,-2.0) -- (7.4,3.4) node[right=1.5mm] {$\mathbb{R}$};
  \node[ponto] (xomega) at (7.4,0.4) {};
  \node[left=2mm] at (7.4,0.1) {$X(\omega)$};

  \draw[->, line width=0.5pt] (12.6,-2.0) -- (12.6,3.4) node[right=1.5mm] {$\mathbb{R}$};
  \node[ponto] (gxomega) at (12.6,0.4) {};
  \node[right=2mm] at (12.6,0.4) {$g(X(\omega))$};

  \draw[->, line width=0.5pt, shorten <=0.3cm, shorten >=0.3cm]
    (omega) to[out=20, in=160] node[pos=0.5, above=1mm] {$X$} (xomega);
  \draw[->, line width=0.5pt, shorten <=0.3cm, shorten >=0.3cm]
    (xomega) to[out=20, in=160] node[pos=0.5, above=1mm] {$g$} (gxomega);
\end{tikzpicture}
\caption{Composição $Y=g\circ X$.}
\label{fig:composicao-y-gx}
\end{figure}


Antes de responder esta pergunta, vamos recordar a definição de
variável aleatória no caso geral.

\begin{definition}[Variável aleatória]
\label{def:variavel-aleatoria}
Seja $(\Omega,\mathcal{F},\mathbb{P})$ um espaço de probabilidade. Dizemos
que $X$ é uma \emph{variável aleatória} se $X$ é uma função
$X:\Omega\to\mathbb{R}$ tal que
\begin{equation}
\label{eq:variavel-aleatoria}
  \{\omega\in\Omega : X(\omega)\leqslant x\}\in\mathcal{F}
  \quad\text{para todo } x\in\mathbb{R}.
\end{equation}
\end{definition}

Segundo a definição acima para que a composição $Y = g(X)$ seja uma variável aleatória 
devemos mostrar que $Y$ satisfaz a condição \eqref{eq:variavel-aleatoria}. 
\begin{figure}[H]
\centering
\begin{tikzpicture}[
  >=Stealth,
  ponto/.style={circle, fill=black, inner sep=0pt, minimum size=1.2mm},
  blob/.style={use Hobby shortcut, closed, line width=0.6pt},
]
  % Faixa do intervalo (-infty,x] na segunda reta
  \def\highlightwidth{8pt}
  \pgfmathsetlengthmacro{\highlightradius}{0.5*\highlightwidth}
  \fill[gray!40, opacity=0.44]
    ([xshift=-\highlightradius]11.8,1.2)
      -- ([xshift=\highlightradius]11.8,1.2)
      -- ([xshift=\highlightradius]11.8,-2.0)
      arc (0:-180:\highlightradius)
      -- cycle;

  % Espaço amostral Omega
  \draw[blob]
    (-1.58,1.2) .. (-0.11,-1.25) .. (1.04,-0.62) .. (2.53,-1.42) .. (3.61,0.13)
    .. (3.88,1.41) .. (3.2,2.12) .. (1.72,2.36) .. (0.23,2.55);
  \node at (0.6,3.0) {$\Omega$};

  % Pré-imagem da composta (região de Omega)
  \filldraw[gray!40, opacity=0.60, blob]
    (0.96,-0.19) .. (1.32,0.51) .. (2.04,0.39) .. (2.64,0.75)
    .. (2.58,1.47) .. (1.98,1.77) .. (1.2,1.7);
  \node at (1.5,1.1) {\small $\{g(X)\leqslant x\}$};

  % Primeira reta
  \draw[-stealth, line width=0.5pt] (7.0,-2.0) -- (7.0,3.4) node[right=1.5mm] {$\mathbb{R}$};
  
  % Segunda reta
  \draw[-stealth, line width=0.5pt] (11.8,-2.0) -- (11.8,3.4) node[right=1.5mm] {$\mathbb{R}$};
  \fill (11.8,1.2) circle (1.3pt);
  \node[left=2mm] at (11.8,1.2) {$x$};

  
  \draw[-stealth] (3.7,2.2) to [out=30, in=150] node[midway, above] {$X$} (6.5,2) ;
  \draw[-stealth] (7.5,2.2) to [out=30, in=150] node[midway, above] {$g$} (11.5,2.2) ;

  % Seta da composta
  \draw[-stealth] (2.5,0.8) to [out=-30, in=-150]  (11.4,-0.3);
  \node at (9,-0.45) {$Y=g\circ X$};  
\end{tikzpicture}
\caption{Pré-imagem de $(-\infty,x]$ pela composta $Y=g\circ X$.}
\label{fig:pre-imagem-composta}
\end{figure}


Note que a única exigência feita sobre $X$ é que as pré-imagens dos 
intervalos $(-\infty,x]$ sejam
eventos. Para que $Y$ herde essa propriedade de $X$, é natural exigir que a função 
$g$ possua a seguinte propriedade: 
a imagem inversa por $X$ de conjuntos da forma $g^{-1}( (-\infty,x] )$
devem ser elementos da $\sigma$-álgebra $\mathcal{F}$.
Vamos mostrar que a classe de funções $g:\mathbb{R}\to\mathbb{R}$ 
que resolvem este problema são as chamadas funções Borel-mensuráveis.
A seguir, apresentamos a definição desta classe de funções. 

\bigskip 

Denotamos por $\mathscr{B}(\mathbb{R})$ a \emph{$\sigma$-álgebra de Borel} de
$\mathbb{R}$, isto é, a menor $\sigma$-álgebra de subconjuntos de
$\mathbb{R}$ que contém todos os intervalos abertos. Seus elementos são
chamados de \emph{borelianos}. Observe que a menor $\sigma$-álgebra 
contendo todos os intervalos da forma $(-\infty,x]$, com
$x\in\mathbb{R}$ coincide com $\mathscr{B}(\mathbb{R})$, pois todo intervalo
aberto pode ser obtido por complementos e uniões enumeráveis de intervalos da forma
$(-\infty,x]$.

\bigskip 

\begin{definition}[Função Borel-mensurável]
\label{def:borel-mensuravel}
Uma função $g:\mathbb{R}\to\mathbb{R}$ é dita \emph{Borel-mensurável} quando
$g^{-1}(B)\in\mathscr{B}(\mathbb{R})$ para todo $B\in\mathscr{B}(\mathbb{R})$.
Equivalentemente, $g$ é Borel-mensurável se, e somente se,
$g^{-1}\bigl((-\infty,x]\bigr)\in\mathscr{B}(\mathbb{R})$ para todo
$x\in\mathbb{R}$.
\end{definition}

\bigskip 
A equivalência registrada na definição é útil na prática: para verificar que
$g$ é Borel-mensurável não é preciso testar todos os borelianos, mas apenas os
intervalos $(-\infty,x]$. A prova desta equivalência segue do mesmo argumento de $\sigma$-álgebra
empregado no lema a seguir. Antes de compor $g$ com $X$, convém registrar que
a condição \eqref{eq:variavel-aleatoria}, exigida apenas para os intervalos
$(-\infty,x]$, vale na verdade para todo boreliano.

\bigskip 

\begin{lemma}\label{lem:pre-imagem-boreliano}
Seja $X$ uma variável aleatória em $(\Omega,\mathcal{F},\mathbb{P})$. Então
$X^{-1}(B)\in\mathcal{F}$ para todo $B\in\mathscr{B}(\mathbb{R})$. Em outras
palavras, $\{X\in B\}\equiv X^{-1}(B)$ é um evento para todo boreliano
$B\subseteq\mathbb{R}$.
\end{lemma}

\begin{proof}
Considere a coleção
\[
  \mathcal{A}_X=\{B\subseteq\mathbb{R}:X^{-1}(B)\in\mathcal{F}\}.
\]
Vamos mostrar que $\mathcal{A}_X$ é uma $\sigma$-álgebra de subconjuntos de
$\mathbb{R}$. De fato, $X^{-1}(\mathbb{R})=\Omega\in\mathcal{F}$, logo
$\mathbb{R}\in\mathcal{A}_X$. Observe também que se $B\in\mathcal{A}_X$, então
\[
  X^{-1}(\mathbb{R}\setminus B)
  =\Omega\setminus X^{-1}(B)\in\mathcal{F},
\]
e portanto $\mathbb{R}\setminus B\in\mathcal{A}_X$, mostrando que a coleção
$\mathcal{A}_X$ é fechada para complementação. Por último, se $(B_n)_{n\geqslant1}$ é
uma sequência arbitrária em $\mathcal{A}_X$, então
\[
  X^{-1}\!\left(\bigcup_{n=1}^{\infty}B_n\right)
  =\bigcup_{n=1}^{\infty}X^{-1}(B_n)\in\mathcal{F},
\]
logo $\bigcup_{n=1}^{\infty}B_n\in\mathcal{A}_X$, o que mostra que a coleção 
$\mathcal{A}_X$ é fechada para uniões enumeráveis.


\bigskip 
Além do mais, segue da definição de variável
aleatória que $\mathcal{A}_X$ contém cada um dos intervalos da forma $(-\infty,x]$, para todo
$x\in\mathbb{R}$. Como esses intervalos geram $\mathscr{B}(\mathbb{R})$,
segue que $\mathscr{B}(\mathbb{R})\subseteq\mathcal{A}_X$, que é a afirmação
do lema.
\end{proof}

\begin{proposition}[Composição com uma função Borel-mensurável]
\label{prop:composicao-borel}
Se $X$ é uma variável aleatória em $(\Omega,\mathcal{F},\mathbb{P})$ e
$g:\mathbb{R}\to\mathbb{R}$ é Borel-mensurável, então $Y=g(X)$ é uma variável
aleatória em $(\Omega,\mathcal{F},\mathbb{P})$.
\end{proposition}

\begin{proof}
Como $Y$ assume valores reais, basta verificar a condição
\eqref{eq:variavel-aleatoria}. Para cada $x\in\mathbb{R}$,
\[
\begin{aligned}
\{Y\leqslant x\}
&=\{\omega\in\Omega:g(X(\omega))\leqslant x\}
\\[0.3cm]
&=\left\{\omega\in\Omega:X(\omega)\in g^{-1}\bigl((-\infty,x]\bigr)\right\}
\\[0.3cm]
&=X^{-1}\Bigl(g^{-1}\bigl((-\infty,x]\bigr)\Bigr).
\end{aligned}
\]
Pela \Cref{def:borel-mensuravel}, o conjunto
$g^{-1}\bigl((-\infty,x]\bigr)$ é um boreliano; pela
\Cref{lem:pre-imagem-boreliano}, sua pré-imagem por $X$ pertence a
$\mathcal{F}$. Logo $\{Y\leqslant x\}\in\mathcal{F}$ para todo
$x\in\mathbb{R}$, isto é, $Y$ é uma variável aleatória.
\end{proof}

Em particular, toda função contínua é Borel-mensurável, pois a pré-imagem de
um aberto é aberta, e o mesmo vale para funções monótonas. Assim, a
\Cref{prop:composicao-borel} cobre os casos usuais. Por causa desta observação daqui em diante
vamos supor sempre que todas as quantidades da forma $Y=g(X)$ 
são variáveis aleatórias. 

\medskip 

A questão principal agora é a
seguinte: se conhecemos a distribuição de $X$, como encontramos a
distribuição de $Y = g(X)$? Se $X$ é discreta com função de probabilidade
$p_X$, então a fórmula apresentada nas notas sobre variáveis aleatórias
discretas fornece a resposta. A seguir, consideramos o caso em que $X$ é
uma variável aleatória absolutamente contínua, com função densidade $f_X$. 
Começamos com um exemplo.

\begin{example}[Transformação afim]\label{ex:transformacao-afim}
  Se $X$ é contínua com função densidade $f_X$, e $g(x) = ax + b$ com
  $a > 0$, então $Y = g(X) = aX + b$ tem função de distribuição dada por
  \[
    \begin{aligned}
      \mathbb{P}(Y \leqslant y) &= \mathbb{P}(aX + b \leqslant y) \\
      &= \mathbb{P}\bigl(X \leqslant a^{-1}(y - b)\bigr) \\
      &= F_X\bigl(a^{-1}(y - b)\bigr).
    \end{aligned}
  \]
  Derivando em relação a $y$, obtemos
  \[
    f_Y(y) = a^{-1} f_X\bigl(a^{-1}(y - b)\bigr) \qquad \text{para } y \in
    \mathbb{R}.
  \]
\end{example}

O próximo teorema generaliza o resultado deste exemplo.

\begin{theorem}[Densidade de uma transformação estritamente crescente]\label{thm:transformacao-crescente}
  Seja $X$ uma variável aleatória contínua com função densidade $f_X$ e seja
  $g : \mathbb{R} \to \mathbb{R}$ estritamente crescente, diferenciável e com
  $g'(x) \neq 0$ para todo $x \in \mathbb{R}$. Então $Y = g(X)$ tem função
  densidade dada por
  \begin{equation}\label{eq:transformacao-crescente}
    f_Y(y) = f_X\bigl(g^{-1}(y)\bigr) \frac{d}{dy}[g^{-1}(y)]
    \qquad \text{para } y \in g(\mathbb{R}),
  \end{equation}
  e $f_Y(y) = 0$ para $y \notin g(\mathbb{R})$. Aqui $g^{-1}$ denota a função
  inversa de $g$, definida em $g(\mathbb{R})$.
\end{theorem}

\begin{proof}
  Primeiro, determinamos a função de distribuição de $Y$:
  \[
    \begin{aligned}
      \mathbb{P}(Y \leqslant y) &= \mathbb{P}(g(X) \leqslant y) \\
      &= \mathbb{P}\bigl(X \leqslant g^{-1}(y)\bigr) \quad \text{pois } g
        \text{ é crescente.}
    \end{aligned}
  \]
  Derivamos essa expressão em relação a $y$ e usamos a regra da cadeia. Como
  $g'(x) \neq 0$, a inversa $g^{-1}$ é diferenciável em $g(\mathbb{R})$ e
  $F_X' = f_X$ nos pontos de diferenciabilidade, de modo que obtemos
  \eqref{eq:transformacao-crescente}.
\end{proof}

Se, no teorema acima, $g$ fosse estritamente decrescente e satisfizesse as
demais hipóteses, o mesmo argumento forneceria que $Y = g(X)$ tem função
densidade
\begin{equation}\label{eq:transformacao-decrescente}
  f_Y(y) = -f_X\bigl(g^{-1}(y)\bigr) \frac{d}{dy}[g^{-1}(y)] \quad \text{para
  } y \in g(\mathbb{R}),
\end{equation}
sendo $f_Y(y) = 0$ para $y \notin g(\mathbb{R})$.

As fórmulas acima dependem fortemente da monotonicidade de $g$. Os demais
casos são tratados de modo mais adequado individualmente, e há, na verdade,
muitas razões para usar sempre o método do próximo exemplo, em vez de recorrer
aos resultados gerais que acabamos de estabelecer.

\begin{example}[Transformação quadrática]\label{ex:transformacao-quadratica}
Suponha que $X$ tenha função densidade $f_X$ e considere
$g(x) = x^2$. Então $Y = g(X) = X^2$ tem função de distribuição
\[
\begin{aligned}
\mathbb{P}(Y \leqslant y) &= \mathbb{P}(X^2 \leqslant y) \\
&= 
\begin{cases} 0 & \text{se } y < 0, \\[0.3cm] 
\mathbb{P}(-\sqrt{y} \leqslant X \leqslant \sqrt{y}) & \text{se } y \geqslant 0.
\end{cases}
\end{aligned}
\]
Consequentemente, $f_Y(y) = 0$ se $y \leqslant 0$, enquanto para
$y > 0$, se essa derivada existir,
\[
\begin{aligned}
f_Y(y) 
= \frac{d}{dy}\mathbb{P}(Y \leqslant y) 
&= \frac{d}{dy}\bigl[F_X(\sqrt{y}) - F_X(-\sqrt{y})\bigr] 
\\[0.3cm]
&= \frac{1}{2\sqrt{y}}\bigl[f_X(\sqrt{y}) + f_X(-\sqrt{y})\bigr].
\end{aligned}
\]
\end{example}

\begin{exercise}\label{exercise:densidades-exponencial}
  Seja $X$ uma variável aleatória com distribuição exponencial de
  parâmetro $\lambda$. Encontre a função densidade de
  \begin{enumerate}[label=(\alph*)]
    \item $A = 2X + 5$,
    \item $B = \exp(X)$,
    \item $C = (1 + X)^{-1}$,
    \item $D = (1 + X)^{-2}$.
  \end{enumerate}
\end{exercise}

\begin{exercise}\label{exercise:normal-qui-quadrado}
  Mostre que se $X$ tem distribuição normal com parâmetros $0$ e
  $1$, então $Y = X^2$ tem distribuição $\chi^2$ com um grau de
  liberdade.
\end{exercise}

\section{Esperança de Variáveis Aleatórias Contínuas}\label{sec:esperanca-va-continuas}

Suponha que uma barra metálica unidimensional tenha densidade $\rho(x)$ no
ponto $x$. Então sua massa é $m = \int \rho(x)\,dx$ e seu centro de
gravidade está na posição $m^{-1} \int x\rho(x)\,dx$. Isso nos conduz
naturalmente à ideia de esperança de uma variável aleatória contínua.

\begin{definition}[Esperança]\label{def:esperanca-continua}
Se $X$ é uma variável aleatória contínua com função de densidade $f_X$,
a \emph{esperança} de $X$ é denotada por $\mathbb{E}[X]$ e definida por
\begin{equation}\label{eq:def-esperanca-continua}
\mathbb{E}[X] = \int_{-\infty}^{\infty} x f_X(x)\,dx,
\end{equation}
sempre que essa integral converge absolutamente, isto é, sempre que
$$\displaystyle\int_{-\infty}^{\infty}|x f_X(x)|\,dx<+\infty.$$
\end{definition}

Assim como no caso das variáveis discretas, a esperança de $X$ é muitas
vezes chamada de \emph{valor esperado} ou \emph{média} de $X$.

Se $X$ é uma variável contínua e $g : \mathbb{R} \to \mathbb{R}$, então
$Y = g(X)$ também é uma variável aleatória, desde que $g$ seja
suficientemente bem comportada. Pode ser difícil calcular $\mathbb{E}[Y]$ a
partir dos primeiros princípios, até porque $Y$ pode não ser nem discreta
nem contínua, de modo que nem a fórmula da esperança no caso discreto nem a
definição acima se aplicam. É de grande utilidade, aqui, o resultado a seguir,
que nos permite calcular $\mathbb{E}[Y]$ diretamente a partir do
conhecimento de $f_X$ e de $g$.

\begin{theorem}[Teorema da Mudança de Variáveis]\label{thm:lotus-continuo}
  Seja $X$ uma variável aleatória contínua com função de densidade $f_X$ e
  seja $g : \mathbb{R} \to \mathbb{R}$ mensurável. Então
  \begin{equation}\label{eq:lotus-continuo}
    \mathbb{E}[g(X)] = \int_{-\infty}^{\infty} g(x) f_X(x)\,dx,
  \end{equation}
  sempre que $\displaystyle\int_{-\infty}^{\infty}|g(x)| f_X(x)\,dx<+\infty$.
\end{theorem}

Este resultado também é conhecido como \emph{lei do estatístico subconsciente}.

\begin{proof}[Esboço de prova]
  O teorema não é difícil de demonstrar, mas a demonstração completa foge ao
  escopo destas notas. Consideramos mais importante entender \emph{por que} ele
  é verdadeiro do que acompanhar uma prova formal. Recorde o teorema da mudança
  de variáveis no caso discreto, visto nas notas sobre esperança de variáveis
  aleatórias discretas: se $Y$ é uma variável aleatória discreta e
  $g : \mathbb{R} \to \mathbb{R}$, então
  \[
    \mathbb{E}[g(Y)] = \sum_x g(x) \mathbb{P}(Y = x).
  \]
  Lembre agora da analogia entre funções de probabilidade e funções de
  densidade: na igualdade acima, substitua $\mathbb{P}(Y = x)$ por
  $f_X(x)\,dx$ e a soma pela integral, obtendo a fórmula do Teorema da
  Mudança de Variáveis. Essa é a razão pela qual o resultado é verdadeiro.
\end{proof}

Assim como no caso das variáveis aleatórias discretas, a média $\mathbb{E}[X]$ de
uma variável aleatória contínua $X$ é um indicativo do ``centro'' da
distribuição de $X$. Como medida do grau de dispersão de $X$ em torno dessa
média, tomamos a \emph{variância} de $X$. Como no caso discreto, exigimos
$\mathbb{E}[X^2]<+\infty$ para que a variância esteja bem definida.

\begin{definition}[Variância]\label{def:variancia-continua}
  A \emph{variância} de uma variável aleatória contínua $X$ tal que
  $\mathbb{E}[X^2]<+\infty$ é definida por
  \begin{equation}\label{eq:def-variancia-continua}
    \operatorname{Var}(X) = \mathbb{E}\big[(X - \mathbb{E}[X])^2\big].
  \end{equation}
\end{definition}

Pelo Teorema da Mudança de Variáveis, \Cref{thm:lotus-continuo}, e observando
que $g(x) = (x - \mu)^2$ é não negativa, obtemos
\[
  \operatorname{Var}(X) = \int_{-\infty}^{\infty} (x - \mu)^2 f_X(x)\,dx,
\]
onde $\mu = \mathbb{E}[X]$. A hipótese $\mathbb{E}[X^2]<+\infty$ garante que
essa integral converge, já que $(x - \mu)^2 \leqslant 2(x^2 + \mu^2)$. Portanto,
\[
\begin{aligned}
\operatorname{Var}(X) &= \int_{-\infty}^{\infty} (x^2 - 2\mu x + \mu^2) f_X(x)\,dx 
\\[0.3cm]
&= \int_{-\infty}^{\infty} x^2 f_X(x)\,dx - 2\mu \int_{-\infty}^{\infty} x f_X(x)\,dx + \mu^2 \int_{-\infty}^{\infty} f_X(x)\,dx 
\\[0.3cm]
&= \int_{-\infty}^{\infty} x^2 f_X(x)\,dx - \mu^2,
\end{aligned}
\]
pois a normalização da densidade nos garante que $\int_{-\infty}^{\infty} f_X(x)\,dx = 1$ e, pela definição de $\mu$, que $\int_{-\infty}^{\infty} x f_X(x)\,dx = \mu$. Assim, obtemos a fórmula usual
\begin{equation}\label{eq:variancia-formula}
  \operatorname{Var}(X) = \mathbb{E}[X^2] - \mathbb{E}[X]^2.
\end{equation}

\begin{example}[Esperança na distribuição uniforme]\label{ex:esperanca-uniforme}
Se $X$ tem a \emph{distribuição uniforme} em $(a, b)$, 
então a sua densidade é $f_X(x) = \frac{1}{b - a}$ para 
$x \in (a, b)$ e $f_X(x) = 0$ caso contrário. Pela definição de esperança,
\[
\begin{aligned}
\mathbb{E}[X] &= \int_{-\infty}^{\infty} x f_X(x)\,dx 
\\[0.3cm]
&= \int_a^b x \frac{1}{b - a}\,dx = \frac{1}{2}(a + b).
\end{aligned}
\]
Como exemplo de uma função de uma variável aleatória, tomemos $Y = \sin X$. Novamente pelo Teorema da Mudança de Variáveis,
\[
\begin{aligned}
\mathbb{E}[Y] &= \int_{-\infty}^{\infty} \sin x \, f_X(x)\,dx 
\\[0.3cm]
&= \int_a^b \frac{\sin x}{b - a}\,dx = \frac{\cos a - \cos b}{b - a}.
\end{aligned}
\]
\end{example}

\begin{example}[Esperança e variância da distribuição exponencial]\label{ex:esperanca-exponencial}
  Se $X$ tem a \emph{distribuição exponencial} de parâmetro $\lambda > 0$, então $f_X(x) = \lambda e^{-\lambda x}$ para $x \geqslant 0$ e $f_X(x) = 0$ para $x < 0$. Pelo Teorema da Mudança de Variáveis,
  \[
    \mathbb{E}[X] = \int_0^{\infty} x \lambda e^{-\lambda x}\,dx
    \qquad\text{e}\qquad
    \mathbb{E}[X^2] = \int_0^{\infty} x^2 \lambda e^{-\lambda x}\,dx.
  \]
Vamos calcular cada uma dessas integrais separadamente, aplicando integração por partes.

Primeiro, vamos calcular o valor médio de $X$. Para a integração por partes
tomamos $u = x$ e $dv = \lambda e^{-\lambda x}\,dx$,
de modo que $du = dx$ e uma primitiva de $dv$ é $v = -e^{-\lambda x}$.
Desta forma temos, pela fórmula de integração por partes que
\[
\begin{aligned}
\int_0^{\infty} x \lambda e^{-\lambda x}\,dx
&= -x e^{-\lambda x}\Bigr|_0^{\infty}
+ \int_0^{\infty} e^{-\lambda x}\,dx 
\\[0.3cm]
&= -x e^{-\lambda x}\Bigr|_0^{\infty}
- \frac{1}{\lambda} e^{-\lambda x}\Bigr|_0^{\infty}
= \frac{1}{\lambda}.
\end{aligned}
\]
Note que na igualdade acima usamos que o termo de fronteira se anula em $x = 0$ e também tem limite zero quando
$x \to \infty$, já que o decaimento exponencial $e^{-\lambda x}$ domina o
crescimento linear de $x$, de modo que $x e^{-\lambda x} \to 0$. Logo
$\mathbb{E}[X] = 1/\lambda$.

\medskip
Para o cálculo do segundo momento, tomamos $u = x^2$ e
$dv = \lambda e^{-\lambda x}\,dx$. Deste modo temos que $du = 2x\,dx$ e
$v = -e^{-\lambda x}$. Aplicando novamente a fórmula da integração por partes
ficamos com
\[
\int_0^{\infty} x^2 \lambda e^{-\lambda x}\,dx
= -x^2 e^{-\lambda x}\Bigr|_0^{\infty}
+ 2\int_0^{\infty} x e^{-\lambda x}\,dx.
\]
Novamente o termo de fronteira se anula, pois a exponencial decai mais
rápido do que $x^2$ cresce. A integral que sobrou é, a menos do fator
$\lambda$, a mesma que já apareceu no cálculo de $\mathbb{E}[X]$.
Assim, realizando novamente a integração por parte com $u = x$ e $dv = e^{-\lambda x}\,dx$, 
temos $du = dx$ e $v = -\lambda^{-1} e^{-\lambda x}$ e portanto 
\[
\begin{aligned}
\int_0^{\infty} x e^{-\lambda x}\,dx
&= -x \frac{1}{\lambda} e^{-\lambda x}\Bigr|_0^{\infty}
+ \frac{1}{\lambda} \int_0^{\infty} e^{-\lambda x}\,dx 
\\[0.3cm]
&= 0 + \frac{1}{\lambda} \cdot \frac{1}{\lambda}
= \frac{1}{\lambda^2}.
\end{aligned}
\]
Da igualdade acima concluímos que $\mathbb{E}[X^2] = 2/\lambda^2$.

Logo, por \eqref{eq:variancia-formula}, a variância de $X$ é
\[
\begin{aligned}
\operatorname{Var}(X) &= \mathbb{E}\big[(X - \mathbb{E}[X])^2\big] \\
&= \mathbb{E}[X^2] - \mathbb{E}[X]^2 = \frac{2}{\lambda^2} - \frac{1}{\lambda^2} = \frac{1}{\lambda^2}.
\end{aligned}
\]
\end{example}

\begin{example}[Esperança e variância da distribuição normal padrão]\label{ex:esperanca-normal}
  Se $X$ tem \emph{distribuição normal} com parâmetros $\mu = 0$ e
  $\sigma^2 = 1$, então
  \[
    \mathbb{E}[X] = \int_{-\infty}^{\infty} x \frac{1}{\sqrt{2\pi}}
    e^{-\frac{1}{2}x^2}\,dx = 0,
  \]
  por simetria do integrando: a função $x \mapsto x e^{-\frac{1}{2}x^2}$
  é ímpar e o intervalo de integração é simétrico em relação à origem.
  Consequentemente, como $\mathbb{E}[X] = 0$, temos
  \[
    \operatorname{Var}(X) = \mathbb{E}[X^2]
    = \frac{1}{\sqrt{2\pi}} \int_{-\infty}^{\infty} x^2
    e^{-\frac{1}{2}x^2}\,dx.
  \]
  Para calcular essa integral, escrevemos o integrando como
  $x \cdot \bigl(x e^{-\frac{1}{2}x^2}\bigr)$ e aplicamos a fórmula da
  integração por partes com $u = x$ e
  $dv = x e^{-\frac{1}{2}x^2}\,dx$. Nesse caso $du = dx$ e
  $v = -e^{-\frac{1}{2}x^2}$, e
  \[
    \begin{aligned}
      \int_{-\infty}^{\infty} x^2 e^{-\frac{1}{2}x^2}\,dx
      &= -x e^{-\frac{1}{2}x^2}\Bigr|_{-\infty}^{\infty}
        + \int_{-\infty}^{\infty} e^{-\frac{1}{2}x^2}\,dx 
      \\[0.3cm]
      &= \int_{-\infty}^{\infty} e^{-\frac{1}{2}x^2}\,dx
        = \sqrt{2\pi}.
    \end{aligned}
  \]
  O termo de fronteira se anula porque
  $x e^{-\frac{1}{2}x^2} \to 0$ quando $|x| \to \infty$, e a última
  igualdade é a normalização da densidade normal padrão. Portanto,
  \[
    \operatorname{Var}(X) = \frac{1}{\sqrt{2\pi}} \cdot \sqrt{2\pi} = 1.
  \]
  Integrações análogas mostram que a distribuição normal com parâmetros
  $\mu$ e $\sigma^2$ tem média $\mu$ e variância $\sigma^2$, como
  era de se esperar.
\end{example}

\begin{example}\label{ex:esperanca-cauchy}
  Neste exemplo vamos mostrar que se $X$ tem \emph{distribuição de Cauchy}, então
  \[
    \mathbb{E}[X] = \int_{-\infty}^{\infty} x \frac{1}{\pi(1 + x^2)}\,dx,
  \]
  não converge, ou seja, uma variável aleatória com distribuição de Cauchy
  não possui valor médio bem-definido. 
  
  \medskip 
  De fato, para cada $M, N > 0$ fixados, temos
  \[
    \begin{aligned}
      \int_{-M}^{N} x \frac{1}{\pi(1 + x^2)}\,dx &=
        \left[\frac{1}{2\pi}\log(1 + x^2)\right]_{-M}^{N} \\
      &= \frac{1}{2\pi}\log\frac{1 + N^2}{1 + M^2} = I(M, N).
    \end{aligned}
  \]
  Portanto o limite de $I(M, N)$ quando $M, N \to \infty$ depende da
  maneira pela qual $M$ e $N$ tendem a $+\infty$. Se fizermos primeiro
  $M \to \infty$ e depois $N \to \infty$, então $I(M, N) \to -\infty$;
  se a passagem ao limite for feita na ordem inversa, então
  $I(M, N) \to +\infty$. Portanto, \emph{a distribuição de Cauchy não possui
  esperança bem-definida}. Em contrapartida, existem muitas funções de $X$ com esperança
  finita. Por exemplo, se $Y = \arctan X$, então
  \[
    \begin{aligned}
      \mathbb{E}[Y] &= \int_{-\infty}^{\infty} \arctan x \frac{1}{\pi(1 +
        x^2)}\,dx \\
      &= \int_{-\frac{1}{2}\pi}^{\frac{1}{2}\pi} \frac{v}{\pi}\,dv, 
        \quad \text{onde } v = \arctan x \\
      &= 0.
    \end{aligned}
  \]
\end{example}

\section{Uma aplicação: A Desigualdade de Cauchy-Schwarz}\label{sec:cauchy-schwarz}

A identidade \eqref{eq:variancia-formula} e o fato de que a variância é
sempre não negativa permitem deduzir um resultado clássico da análise. A
dedução passa pela forma probabilística da desigualdade, enunciada a
seguir.

\begin{proposition}[Cauchy-Schwarz, forma
  probabilística]\label{prop:cauchy-schwarz}
  Seja $X$ uma variável aleatória contínua com função de densidade $f_X$ e
  sejam $g, h : \mathbb{R} \to \mathbb{R}$ funções mensuráveis tais que
  $\mathbb{E}[g(X)^2] < +\infty$ e $\mathbb{E}[h(X)^2] < +\infty$. Então o
  produto $g(X)h(X)$ é integrável e
  \[
    \bigl(\mathbb{E}[g(X)h(X)]\bigr)^2
    \leqslant
    \mathbb{E}[g(X)^2]\,\mathbb{E}[h(X)^2].
  \]
\end{proposition}

\begin{proof}
Começamos por verificar que as esperanças do enunciado fazem sentido.
Ponto a ponto vale $2\,|gh| \leqslant g^2 + h^2$ e, portanto,
\[
\begin{aligned}
2\int_{-\infty}^{\infty} |g(x)h(x)| f_X(x)\,dx
&\leqslant \int_{-\infty}^{\infty} g(x)^2 f_X(x)\,dx
+ \int_{-\infty}^{\infty} h(x)^2 f_X(x)\,dx
\\[0.3cm]
&= \mathbb{E}[g(X)^2] + \mathbb{E}[h(X)^2] < +\infty.
\end{aligned}
\]
Assim a integral de $g(x)h(x)f_X(x)$ converge absolutamente e, pelo
\Cref{thm:lotus-continuo}, o produto $g(X)h(X)$ é integrável.

\medskip
Agora usamos o fato de que a variância é sempre não negativa. Fixe
$\lambda \in \mathbb{R}$ e considere a variável aleatória 
$Z = g(X) - \lambda\,h(X)$. Aplicando a desigualdade
$(a - b)^2 \leqslant 2a^2 + 2b^2$ a variável aleatória $Z$ concluímos que
$Z^2 \leqslant 2g(X)^2 + 2\lambda^2 h(X)^2$. Logo, segue diretamente da hipótese que
$\mathbb{E}[Z^2] < +\infty$ e a variância de $Z$ está bem definida.
Pela fórmula da variância \eqref{eq:variancia-formula} e pelo fato que
$\operatorname{Var}(Z) \geqslant 0$, temos
\[
\mathbb{E}[Z^2]
= \operatorname{Var}(Z) + \mathbb{E}[Z]^2 \geqslant 0.
\]
Expandindo o quadrado de $Z$ e usando a linearidade da integral,
obtemos, para cada $\lambda \in \mathbb{R}$,
\[
\mathbb{E}[Z^2]
= \mathbb{E}[g(X)^2] - 2\lambda\,\mathbb{E}[g(X)h(X)]
+ \lambda^2\,\mathbb{E}[h(X)^2] \geqslant 0.
\]
Ou seja, o polinômio do segundo grau em $\lambda$ dado pelo lado direito
acima é não negativo em toda a reta.

\medskip
Para concluir, consideramos dois casos. Se $\mathbb{E}[h(X)^2] = 0$,
tomemos $\varepsilon > 0$. Como a desigualdade
\[
0 \leqslant 
\left( \sqrt{\varepsilon}\,|g| -\frac{1}{\sqrt{\varepsilon}}|h|\right)^2
= \varepsilon\, g^2 - 2\,|gh| + \frac{1}{\varepsilon} h^2,
\]
vale em cada ponto, podemos afirmar que 
\[ 
2\,|gh| \leqslant \varepsilon\, g^2 + \frac{1}{\varepsilon} h^2.
\] 
Usando a desigualdade acima temos a seguinte estimativa
\[
2\big|\mathbb{E}[g(X)h(X)]\big|
\leqslant 2\,\mathbb{E}\big[|g(X)h(X)|\big]
\leqslant \varepsilon\,\mathbb{E}[g(X)^2]
+ \frac{1}{\varepsilon}\mathbb{E}[h(X)^2]
= \varepsilon\,\mathbb{E}[g(X)^2],
\]
onde na última igualdade usamos que $\mathbb{E}[h(X)^2]$ se anula.
Fazendo $\varepsilon \to 0$, concluímos que $\mathbb{E}[g(X)h(X)]$
também se anula. Neste caso concluímos que a desigualdade do enunciado vale como
igualdade. 

Por outro lado, se $\mathbb{E}[h(X)^2] > 0$, então o polinômio acima tem
coeficiente dominante positivo e também é não-negativo em toda a reta. Logo o
seu discriminante é não-positivo, ou seja,
\[
4\bigl(\mathbb{E}[g(X)h(X)]\bigr)^2
- 4\,\mathbb{E}[g(X)^2]\,\mathbb{E}[h(X)^2]
\leqslant 0.
\]
Dividindo por $4$ chega-se à desigualdade do enunciado.
\end{proof}

A versão clássica desta desigualdade, pode ser obtida considerando
o caso particular em que $X$ é uma variável aleatória uniforme.

\begin{corollary}[Cauchy-Schwarz para
  integrais]\label{cor:cauchy-schwarz-integral}
  Sejam $a < b$ e $g, h : [a, b] \to \mathbb{R}$ funções tais que $g^2$ e
  $h^2$ são integráveis em $[a, b]$. Então $gh$ é integrável em $[a, b]$ e
  \[
    \biggl(\int_a^b g(x)h(x)\,dx\biggr)^{\!2}
    \leqslant
    \int_a^b g(x)^2\,dx \int_a^b h(x)^2\,dx.
  \]
\end{corollary}

\begin{proof}
  Seja $X$ com distribuição uniforme em $(a, b)$, como no
  \Cref{ex:esperanca-uniforme}. Identificando $g$ e $h$ com suas extensões
  por zero fora de $[a, b]$, as hipóteses garantem
  $\mathbb{E}[g(X)^2] < +\infty$ e $\mathbb{E}[h(X)^2] < +\infty$. Pela
  parte inicial da \Cref{prop:cauchy-schwarz}, o produto $g(X)h(X)$ é
  integrável, isto é, $gh$ é integrável em $[a, b]$. Pelo
  \Cref{thm:lotus-continuo}, valem então as três igualdades a seguir.
  \[
    \begin{aligned}
      \mathbb{E}[g(X)h(X)] &= \frac{1}{b - a}\int_a^b g(x)h(x)\,dx,
      \\[0.3cm]
      \mathbb{E}[g(X)^2] &= \frac{1}{b - a}\int_a^b g(x)^2\,dx,
      \\[0.3cm]
      \mathbb{E}[h(X)^2] &= \frac{1}{b - a}\int_a^b h(x)^2\,dx.
    \end{aligned}
  \]
  Substituindo estas três igualdades na \Cref{prop:cauchy-schwarz} e
  simplificando os fatores $(b - a)^{-2}$, obtemos a desigualdade do
  enunciado.
\end{proof}

\begin{remark}\label{rem:cauchy-schwarz-moral}
  O caso $h \equiv 1$ da \Cref{prop:cauchy-schwarz} é exatamente a
  afirmação de que a variância de $g(X)$ é não-negativa. De fato,
  \[
    \mathbb{E}[g(X)]^2 \leqslant \mathbb{E}[g(X)^2],
  \]
  e portanto segue da fórmula \eqref{eq:variancia-formula} que
  $\operatorname{Var}(g(X)) \geqslant 0$. 
\end{remark}

\section{Exercícios propostos}\label{sec:exercicios-propostos}

\begin{exercise}\label{exercise:esperanca-arcoseno}
  Mostre que uma variável aleatória cuja função de densidade é dada por
  \[
    f(x) = \begin{cases}
      \dfrac{1}{\pi\sqrt{x(1 - x)}} & \text{se } 0 < x < 1, \\[2mm]
      0 & \text{caso contrário,}
    \end{cases}
  \]
  tem média igual a $\frac{1}{2}$.
\end{exercise}

\begin{exercise}\label{exercise:densidade-quadratica}
  A variável aleatória $X$ tem função de densidade
  \[
    f(x) = c\,x(1 - x) \quad \text{para } 0 \leqslant x \leqslant 1.
  \]
  Determine a constante $c$ e calcule a média e a variância de $X$.
\end{exercise}

\begin{exercise}\label{exercise:esperanca-exponencial-quadrado}
  Se $X$ tem distribuição normal com média $0$ e variância $1$,
  calcule $\mathbb{E}\!\left[\,e^{2X}\,\right]$.
\end{exercise}

\begin{exercise}
  \label{exercise:transformacao-normal}
  Sejam $X\sim N(0,1)$, $\mu\in\mathbb{R}$ e $0<\sigma<+\infty$
  constantes. Seguindo o roteiro abaixo mostre que 
  $Y\equiv\sigma X+\mu$ é uma variável aleatória
  com distribuição $N(\mu,\sigma^2)$.
  \begin{enumerate}[label=\textbf{(\alph*)}, leftmargin=*, itemsep=0.6em]
    \item Usando a função auxiliar $g(x)=\sigma x+\mu$ argumente que $Y=g(X)$ é uma
      variável aleatória.
    \item Mostre que $g$ é uma função estritamente crescente. Em seguida, determine
      sua imagem $g(\mathbb{R})$, sua inversa $g^{-1}$ e a derivada
      $(g^{-1})'(y)$.
    \item Aplique o \Cref{thm:transformacao-crescente} para escrever a
      densidade $f_Y$ em termos de $f_X$, de $g^{-1}$ e de $(g^{-1})'$.
    \item Usando a expressão explícita da densidade da normal padrão dada no
      \Cref{ex:esperanca-normal} e os itens anteriores mostre que, para todo $y\in\mathbb{R}$, temos
      \[
        f_Y(y)
        =\frac{1}{\sqrt{2\pi\sigma^2}}
          \exp\!\left(-\frac{(y-\mu)^2}{2\sigma^2}\right).
      \]
      e conclua que $Y\sim N(\mu,\sigma^2)$.
    \item \emph Usando o
      \Cref{thm:lotus-continuo} e os valores de $\mathbb{E}[X]$ e
      $\operatorname{Var}(X)$ calculados no \Cref{ex:esperanca-normal},
      mostre que $\mathbb{E}[Y]=\mu$ e $\operatorname{Var}(Y)=\sigma^2$.
  \end{enumerate}
\end{exercise}

\begin{exercise}\label{exercise:sigma-algebra-gerada}
  Recorde que uma coleção $\mathcal{A}$ de subconjuntos de $\mathbb{R}$ é uma
  \emph{$\sigma$-álgebra de subconjuntos de $\mathbb{R}$}, ou simplesmente uma
  \emph{sub-$\sigma$-álgebra de $\mathbb{R}$}, quando
  \begin{enumerate}[label=(\roman*)]
    \item $\mathbb{R} \in \mathcal{A}$;
    \item se $A \in \mathcal{A}$, então $\mathbb{R} \setminus A \in \mathcal{A}$;
    \item se $A_1, A_2, \dots \in \mathcal{A}$, então
      $\bigcup_{n \geqslant 1} A_n \in \mathcal{A}$.
  \end{enumerate}
  No texto, $\mathscr{B}(\mathbb{R})$ foi apresentada como ``a menor
  $\sigma$-álgebra de subconjuntos de $\mathbb{R}$ que contém todos os
  intervalos abertos''. Nosso objetivo é dar sentido preciso a essa descrição e
  mostrar que ela de fato define um único objeto. Siga o roteiro abaixo.

  \begin{enumerate}[label=\textbf{(\alph*)}, leftmargin=*, itemsep=0.6em]
    \item Seja $I \neq
    \emptyset$ um conjunto de índices e suponha que, para cada $i \in I$,
    $\mathcal{A}_i$ seja uma sub-$\sigma$-álgebra de $\mathbb{R}$. Considere a
    coleção
    \[
      \mathcal{A} = \bigcap_{i \in I} \mathcal{A}_i.
    \]
    Mostre que $\mathcal{A}$ é uma sub-$\sigma$-álgebra de $\mathbb{R}$,
    verificando separadamente que:
    \begin{enumerate}[label=(\roman*)]
      \item $\mathbb{R} \in \mathcal{A}$;
      \item se $A \in \mathcal{A}$, então $\mathbb{R} \setminus A \in \mathcal{A}$;
      \item se $A_1, A_2, \dots \in \mathcal{A}$, então
        $\bigcup_{n \geqslant 1} A_n \in \mathcal{A}$.
    \end{enumerate}
    Em cada item, use apenas a definição de interseção e o fato de que cada
    $\mathcal{A}_i$ satisfaz as condições (i)--(iii) acima. A hipótese $I \neq
    \emptyset$ é essencial: explique em qual passo ela é usada.

    \item Seja $\mathcal{C} \subseteq \mathcal{P}(\mathbb{R})$ uma
    coleção arbitrária de subconjuntos de $\mathbb{R}$ e defina a família
    \[
      \mathcal{F}_{\mathcal{C}}
      = \bigl\{\mathcal{A} : \mathcal{A} \text{ é sub-$\sigma$-álgebra de }
        \mathbb{R} \text{ e } \mathcal{C} \subseteq \mathcal{A}\bigr\}.
    \]
    \begin{enumerate}[label=(\roman*)]
      \item Mostre que $\mathcal{F}_{\mathcal{C}} \neq \emptyset$, exibindo um
        elemento explícito (por exemplo, $\mathcal{P}(\mathbb{R})$). Por que
        este passo é indispensável antes de tomar a interseção da família?
      \item Definimos a \emph{$\sigma$-álgebra gerada por $\mathcal{C}$} por
        \[
          \sigma(\mathcal{C})
          = \bigcap_{\mathcal{A} \in \mathcal{F}_{\mathcal{C}}} \mathcal{A}.
        \]
        Conclua dos itens (a) e (b)(i) que $\sigma(\mathcal{C})$ é uma
        sub-$\sigma$-álgebra de $\mathbb{R}$.
      \item Mostre que $\mathcal{C} \subseteq \sigma(\mathcal{C})$.
      \item Seja $\mathcal{A}$ uma
        sub-$\sigma$-álgebra de $\mathbb{R}$ com $\mathcal{C} \subseteq
        \mathcal{A}$. Mostre que $\sigma(\mathcal{C}) \subseteq \mathcal{A}$. Conclua que
        $\sigma(\mathcal{C})$ é a menor sub-$\sigma$-álgebra de $\mathbb{R}$
        que contém $\mathcal{C}$ e que esta propriedade a caracteriza de forma
        única.
    \end{enumerate}

    \item Tome agora para $\mathcal{C}$
    a coleção de todos os intervalos abertos de $\mathbb{R}$:
    \[
      \mathcal{C}_{\text{ab}}
      = \bigl\{(a,b) : a, b \in \mathbb{R},\ a < b\bigr\}.
    \]
    Pelo item (b), a $\sigma$-álgebra $\sigma(\mathcal{C}_{\text{ab}})$ está
    bem definida. Definimos então
    \[
      \mathscr{B}(\mathbb{R}) = \sigma(\mathcal{C}_{\text{ab}}).
    \]
    Interprete esta definição à luz do item (b)(iv): ela é precisamente o
    conteúdo da expressão ``a menor $\sigma$-álgebra de subconjuntos de
    $\mathbb{R}$ que contém todos os intervalos abertos''.

    \item Fixe $a < b$ em $\mathbb{R}$. Mostre que:
    \begin{enumerate}[label=(\roman*)]
      \item $[a,b] = \bigcap_{n \geqslant 1} \bigl(a - \tfrac{1}{n},\,
        b + \tfrac{1}{n}\bigr)$, isto é, um intervalo fechado é uma
        interseção enumerável de intervalos abertos;
      \item $(a,b) = \bigcup_{n \geqslant 1} \bigl[a + \tfrac{1}{n},\,
        b - \tfrac{1}{n}\bigr]$ (descartando os $n$ para os quais
        $1/n \geqslant (b-a)/2$, de modo que os intervalos não sejam
        vazios), isto é, um intervalo aberto é uma união enumerável de
        intervalos fechados;
      \item $[a,b) = \bigcap_{n \geqslant 1} \bigl(a - \tfrac{1}{n},\, b\bigr)$
        e $(a,b] = \bigcap_{n \geqslant 1} \bigl(a,\, b + \tfrac{1}{n}\bigr)$;
        em particular, um conjunto unitário é um boreliano, pois
        $\{a\} = \bigcap_{n \geqslant 1}\bigl(a - \tfrac{1}{n},\,
        a + \tfrac{1}{n}\bigr)$;
      \item generalize (ii): mostre que todo subconjunto aberto
        $U \subseteq \mathbb{R}$ é união enumerável de intervalos fechados.
        \emph{Dica:} use a densidade de $\mathbb{Q}$ para escrever $U$ como
        união enumerável de intervalos abertos com extremos racionais e
        aplique (ii) a cada um deles.
    \end{enumerate}
    Conclua que, no que diz respeito à $\sigma$-álgebra gerada, não faz
    diferença partir de intervalos abertos ou de intervalos fechados: cada
    tipo se obtém do outro por uniões e interseções enumeráveis.

    \item Reunindo os itens anteriores,
    mostre que todas as coleções abaixo geram a mesma $\sigma$-álgebra, a
    saber, $\mathscr{B}(\mathbb{R})$:
    \begin{enumerate}[label=(\roman*)]
      \item os intervalos abertos $\{(a,b) : a < b\}$;
      \item os intervalos fechados $\{[a,b] : a < b\}$;
      \item os intervalos semi-abertos $\{[a,b) : a < b\}$ e, analogamente,
        $\{(a,b] : a < b\}$;
      \item os intervalos da forma $(-\infty, x]$, com $x \in \mathbb{R}$;
      \item os subconjuntos abertos de $\mathbb{R}$.
    \end{enumerate}
    \emph{Dica:} para provar que duas coleções $\mathcal{C}_1$ e
    $\mathcal{C}_2$ geram a mesma $\sigma$-álgebra, é suficiente, pela
    minimalidade do item (b)(iv), mostrar que
    $\mathcal{C}_1 \subseteq \sigma(\mathcal{C}_2)$ e
    $\mathcal{C}_2 \subseteq \sigma(\mathcal{C}_1)$. Use o item (d) para
    obter cada tipo de intervalo a partir dos demais por operações
    enumeráveis. Para o item (iv), escreva cada $(a,b)$ a partir de
    intervalos $(-\infty, x]$ e, reciprocamente, obtenha $(-\infty, x]$ a
    partir de intervalos abertos. Isto justifica a afirmação do texto de que
    a menor $\sigma$-álgebra que contém os intervalos $(-\infty, x]$
    coincide com $\mathscr{B}(\mathbb{R})$.
  \end{enumerate}
\end{exercise}

\bibliographystyle{plain}
\bibliography{../arquivos-aux/referencias}
\end{document}
